\section{Выводы и обсуждение результатов}

В ходе выполнения лабораторной работы были исследованы явления дифракции света на ультразвуковой волне в жидкости. Основные цели работы включали:
\begin{itemize}
    \item Изучение дифракции света на синусоидальной акустической решётке
    \item Наблюдение фазовой решётки методом тёмного поля
\end{itemize}

Все поставленные задачи были успешно выполнены. В результате экспериментов получены значения скорости звука в воде различными методами:

Скорость звука в воде (первый метод):
\begin{itemize}
    \item $v = (1470 \pm 282)$ м/с, \ $\varepsilon = 19$ \% 
\end{itemize}

Скорость звука в воде (второй метод, среднее значение):
\begin{itemize}
    \item $v = (1509 \pm 25)$ м/с, \ $\varepsilon = 1.7$ \% 
\end{itemize}

Скорость звука в воде (метод тёмного поля):
\begin{itemize}
    \item $v = (1470 \pm 20)$ м/с, \ $\varepsilon = 1.4$ \% 
\end{itemize}

Сопоставляя полученные результаты, можно отметить хорошее согласование значений скорости звука в воде, определённых разными методами. Все три метода дали результаты, близкие к теоретическому значению скорости звука в воде при комнатной температуре ($\approx 1.5$ км/с). При этом первый метод имеет существенно большую погрешность, что связано с трудностью точного определения расстояния между дифракционными максимумами.

В эксперименте наблюдался линейный характер зависимости длины ультразвуковой волны от обратной частоты, что подтверждает справедливость соотношения $\Lambda = v/\nu$, где $v$ — скорость звука, $\nu$ — частота ультразвука.

Метод тёмного поля показал свою эффективность для исследования акустических решёток, позволив наблюдать характерную двойную периодичность структуры и получить точные количественные данные о параметрах ультразвуковой волны. Особенностью этого метода является удвоение частоты наблюдаемой структуры, что хорошо согласуется с теоретическими предсказаниями.
